\documentclass[11pt]{article}
\usepackage[margin=1in]{geometry}
\usepackage{amsmath}
\usepackage{amssymb}
\usepackage{booktabs}
\usepackage{enumitem}
\usepackage{float}
\usepackage{graphicx}
\usepackage{tabularx}
\usepackage[hidelinks]{hyperref}
\newcommand{\pra}{\mathrm{PRA}}
\newcommand{\refhandle}{\mathrm{REF}}
\newcommand{\kv}{\mathrm{KV}}
\newcommand{\q}{\mathbf{Q}}
\newcommand{\kmat}{\mathbf{K}}
\newcommand{\vmat}{\mathbf{V}}

\setlength{\emergencystretch}{3em}

\title{Retrofitting Pretrained Transformers with Progressive Retrieval Attention:\
Thin Hugging Face Adapters and a Qwen Correctness Study}
\author{Jorge Simao\\EInnovator R\&D -- Center for the Sciences of Computation, Cognition, and Complexity}
\date{\today}
\hypersetup{pdftitle={Retrofitting Pretrained Transformers with Progressive Retrieval Attention},pdfauthor={Jorge Simao}}

\begin{document}
\maketitle

\begin{abstract}
Progressive Retrieval Attention (PRA) selects sparse, layer-native key/value (K/V) memory
from references or displaced prompt history. Controlled models establish the mechanism but do
not show that it can preserve a real pretrained transformer's attention semantics. We begin a
pretrained integration study with a shared PRA execution core and thin Hugging Face attention
adapters. The first adapter wraps one upper layer of the frozen 596M-parameter
\texttt{Qwen/Qwen3-0.6B} checkpoint. When memory is disabled, the wrapper delegates to the
original module and obtains bit-exact logits, hidden states, greedy generation, and all 28
generation-cache shapes. When memory is enabled, it reuses Qwen's pretrained projections,
Q/K normalization, rotary embedding, grouped-query layout, causal mask, cache update, eager
attention function, and output projection. Stored memory retains eight native K/V heads rather
than expanding to 16 query heads. A CUDA smoke run materializes 22 reference tokens from CPU,
generates finite text, and executes a 408-token logical prompt as a 280-token implicit head plus
a 128-token direct tail with zero configured native-limit violations. These are transport and
compatibility results, not retrieval-quality results. In one HotpotQA and one QASPER example,
mean-gist routing selected no annotated evidence; HotpotQA failed under every tested baseline,
and the QASPER answer appeared under dense context and PRA without establishing causation.
The first milestone therefore supports the adapter boundary and native-KV transport, while
identifying frozen zero-shot routing as the next blocking scientific problem before Llama or
Gemma integration.
\end{abstract}

\section{Introduction}

Pretrained decoder families differ in projection layout, rotary implementation, grouped-query
attention (GQA), masks, sliding windows, and generation-cache conventions. Replacing a complete
model class to add memory duplicates rapidly changing library code and makes parity difficult
to audit. The alternative studied here is narrower:

\begin{center}
normal pretrained model $\rightarrow$ wrap selected attention modules $\rightarrow$
shared PRA core.
\end{center}

The family adapter owns only native attention semantics. The shared core owns references,
chunk gists, exact routing, the materialization budget, selected K/V transfer, and metrics.
This boundary follows the broader transformer interface of queries attending to keys and
values~\cite{vaswani2017attention}, while preserving Qwen's pretrained RoPE and GQA choices
rather than approximating them in a parallel implementation~\cite{su2021roformer,ainslie2023gqa,qwen2025qwen3}.

This paper begins with Qwen because a sub-billion-parameter checkpoint makes exact tests
practical on constrained hardware while still exercising modern RoPE and GQA. It does not yet
claim cross-family portability or useful long-context QA. The contribution is a tested first
adapter and an evidence hierarchy that prevents working transport from being mistaken for
working retrieval.

\section{PRA Recap and Prior Contracts}

For layer $\ell$, PRA stores reference chunks as native
$K^{(\ell)}_c,V^{(\ell)}_c\in\mathbb{R}^{H_{kv}\times T_c\times d_h}$ and one or more compact
gists $g^{(\ell)}_c$. A query representation $r^{(\ell)}$ scores the gists, and a global
budgeter selects whole chunks. Materialized memory precedes the direct K/V sequence under one
softmax:
\begin{equation}
  \operatorname{Attn}(Q,[K_M;K_D],[V_M;V_D])
  = \operatorname{softmax}\!\left(\frac{Q[K_M;K_D]^\top}{\sqrt{d_h}}+M\right)[V_M;V_D].
\end{equation}
This is not a separately normalized cross-attention residual. The native output projection is
applied once to the combined result.

Paper~1 established controlled native-K/V routing and materialization. Paper~1.5 separated
logical source position from contextual fragmentation. Continuous sources use
$p_i=p_{\mathrm{logical\ start}}+i$, and ordinary post-RoPE keys remain the canonical cache
state. The present adapter preserves that contract by capturing each selected HF layer's
actual post-RoPE keys and by assigning source-relative positions across bounded encoding
blocks.

\section{Adapter Architecture}

\begin{figure}[H]
\centering
\fbox{\begin{minipage}{0.92\linewidth}
\centering
\textbf{Logical source / long prompt}\\[2pt]
$\downarrow$ bounded source encoding with logical offsets\\[2pt]
$\downarrow$ family-native projections, Q/K normalization, and RoPE\\[2pt]
$\downarrow$ native $H_{kv}$ K/V + smaller routing chunks + mean gists\\[2pt]
$\downarrow$ shared exact router $\rightarrow$ shared whole-chunk budgeter\\[2pt]
$\downarrow$ selected CPU/GPU native K/V materialization\\[2pt]
$\downarrow$ thin family adapter prepends memory to local/cache K/V\\[2pt]
$\downarrow$ family-native eager attention and pretrained output projection
\end{minipage}}
\caption{Paper~2 boundary. Model-family semantics stay in a thin adapter; references,
routing, budgeting, materialization, residency, and metrics stay in one shared PRA core.}
\label{fig:adapter}
\end{figure}

\subsection{Shared execution core}

The controlled model's former monolithic attention forward pass was divided into reusable
operations. \texttt{prepare\_pra\_query} selects the newest valid query and maps GQA query
groups into the native K/V routing width. \texttt{route\_memory} calls the exact packed-index
router. \texttt{budget\_selected\_memory} enforces
\begin{equation}
 L_{\mathrm{direct}}+L_{\mathrm{materialized}}+L_{\mathrm{reserve}}
 \leq L_{\mathrm{native,max}}.
\end{equation}
\texttt{materialize\_selected\_kv} removes duplicate overlap, preserves row isolation, and
moves only retained detail K/V. \texttt{combine\_local\_and\_memory\_kv} pads variable rows
and prepends an additive all-visible memory mask to the family's existing local causal mask.
The controlled model and HF adapters now call this same core.

\subsection{Thin HF contract}

The abstract adapter exposes six family operations: project Q/K/V, apply native position
encoding, normalize tensor layout, combine the native mask, invoke PRA attention, and project
the output. Disabled memory takes an even narrower path: call the original attention object
unchanged. This design preserves pretrained parameters and avoids checkpoint conversion.

The initial Qwen file has 215 physical lines (193 nonblank, non-comment lines). It introduces
no learned parameter. Its family-specific branches distinguish Qwen2-style projections from
Qwen3's Q/K normalization; routing, budgeting, materialization, residency, and metrics remain
outside the file. This is a portability measurement, not a code-minimization target.

\section{Qwen Native Semantics}

\subsection{Projection and RoPE}

The adapter calls the wrapped module's \texttt{q\_proj}, \texttt{k\_proj}, and
\texttt{v\_proj}. Qwen3's existing \texttt{q\_norm} and \texttt{k\_norm} are applied before
the installed Transformers \texttt{apply\_rotary\_pos\_emb} function. Reference capture saves
post-RoPE $K$ and position-independent $V$. No independent RoPE implementation is introduced.

\subsection{Grouped-query attention}

The tested checkpoint has $H_q=16$, $H_{kv}=8$, and $d_h=128$. Cache objects retain
$[B,8,T,128]$ K/V. For routing only, each pair of query heads sharing one K/V head is averaged,
producing a vector in the same $8\times128$ space as a flattened key gist. K/V expansion occurs
inside Qwen's eager attention function, never in persistent PRA storage. A synthetic replay
test compares this path with explicit K/V head repetition and matches within $10^{-6}$.

\subsection{HF cache and masks}

Prefill and single-token decode both use the wrapped module's HF \texttt{Cache.update} call.
PRA memory is prepended after that update, so direct generation history appears once. A decode
test with four prompt tokens and two generated tokens observes five direct cache positions at
the last invocation, plus four separate memory positions. The original causal mask remains on
local/cache positions; selected historical memory receives a visible prefix mask. Memory-
disabled generation delegates entirely to the original module.

\section{Experimental Setup}

Table~\ref{tab:model} records the first checkpoint exactly. Tests use synthetic Qwen3 configs
offline. The pretrained run uses Python 3.10.11, PyTorch 2.12.1+CUDA~12.6, Transformers 4.55.4,
eager attention, FP16, and an NVIDIA GeForce GTX~950M. The frozen model is not fine-tuned.

\begin{table}[H]
\centering
\caption{Pinned Qwen checkpoint and PRA placement.}
\label{tab:model}
\begin{tabular}{ll}
\toprule
Checkpoint & \texttt{Qwen/Qwen3-0.6B} \\
Revision & \texttt{c1899de289a04d12100db370d81485cdf75e47ca} \\
Parameters / layers & 596,049,920 / 28 \\
Query heads / K/V heads / head width & 16 / 8 / 128 \\
RoPE base / advertised positions & $10^6$ / 40,960 \\
PRA placement & layer 27 only \\
Initial gist / routing policy & mean, one gist/chunk, exact top-$k$ \\
Native limit / direct / selected-memory cap & 256 / 128 / 64 tokens \\
K/V residency & CPU full cache, selected transfer \\
\bottomrule
\end{tabular}
\end{table}

Every run writes JSON and scalar CSV with the Git SHA, model revision, software versions,
device, dtype, layer selection, limits, routing settings, timings, and metrics. The exact
executed commit is retained in each artifact rather than inferred from the manuscript build.

\section{Adapter Parity}

PRA-disabled parity is the hard gate. The baseline is evaluated first; the same in-memory model
is then wrapped, leaving every parameter object in place. Table~\ref{tab:parity} reports exact
results for the five-token prompt ``The capital of Portugal is'' and four-token greedy decode.

\begin{table}[H]
\centering
\caption{Original HF model versus memory-disabled PRA wrapper.}
\label{tab:parity}
\begin{tabular}{lr}
\toprule
Metric & Result \\
\midrule
Maximum / mean absolute logit difference & 0 / 0 \\
Maximum hidden-state difference & 0 \\
Greedy token agreement & 100\% \\
Greedy sequences exactly equal & yes \\
Generation-cache shapes equal & all 28 layers \\
Finite adapted logits & yes \\
\bottomrule
\end{tabular}
\end{table}

Single-run prefill times were 0.447~s before wrapping and 0.136~s after wrapping. This order is
confounded by first-call warm-up and is not a speed comparison. Exact outputs, rather than
latency, are the parity evidence.

\section{Native-KV Reference Smoke}

An explicit 22-token reference about Lisbon is encoded through the frozen model. The selected
layer captures $[1,8,22,128]$ K/V, stores it on CPU, creates a mean gist, and transfers all 22
tokens because the source fits one routing chunk. The physical transfer is 90,112 bytes, exactly
$2\times22\times8\times128\times2$ bytes for FP16 K and V. The run reports 1.6~ms routing,
0.40~ms materialization orchestration, 0.15~ms selected transfer, and 0.11~ms combined attention.
These synchronized phase values are instrumentation checks on one machine, not throughput
claims. Cold reference construction took 0.144~s and the warm query took 0.191~s.

The generated continuation contains ``Lisbon'', but the unmodified model also knows this fact.
Consequently, this observation establishes functional selected-K/V transport and generation,
not a causal retrieval benefit. The exact native-replay claim is instead restricted to the
synthetic tensor test, where GQA memory plus local K/V equals explicit dense composition.

\section{Implicit Head and Bounded Long Prompt}

A repeated prompt tokenizes to 408 logical tokens. The direct cap retains the final 128 tokens;
the first 280 become \texttt{\#\_\_head}. That head is encoded in blocks of at most 64 tokens,
with source-relative positions. The direct tail receives positions 280 through 407. Every
encoding operation and the final selected-memory operation stays under the 256-token configured
native limit; the largest observed source-encoding call is 64 tokens and violations are zero.
The final finite-logit query took 0.258~s. This proves bounded execution and offset propagation
for the frozen Qwen adapter. It does not show that the router selected the information needed by
an arbitrary continuation.

\section{Unrestricted QA Smoke}

We next run one fixed HotpotQA example~\cite{yang2018hotpotqa} and one QASPER
example~\cite{dasigi2021qasper}. Four conditions use the same frozen checkpoint: question-only
truncation, dense text left-truncated to 256 tokens, oracle evidence text-RAG, and PRA with the
question direct and the complete source in native-K/V memory. Oracle text-RAG is an upper
control because it receives annotated evidence. This two-example matrix diagnoses plumbing; it
cannot estimate dataset accuracy.

\begin{table}[H]
\centering
\caption{Unrestricted-QA smoke outcomes. Standard normalized token F1 is shown.}
\label{tab:qa}
\begin{tabularx}{\linewidth}{lXcc}
\toprule
Dataset & Condition & F1 & Annotated evidence selected \\
\midrule
HotpotQA & truncation; dense; oracle text-RAG; PRA & 0; 0; 0; 0 & PRA: 0\% \\
QASPER & truncation; dense; oracle text-RAG; PRA & 0; 0.182; 0; 0.182 & PRA: 0\% \\
\bottomrule
\end{tabularx}
\end{table}

HotpotQA's expected answer is ``yes''; every condition begins with ``No''. QASPER's expected
answer is ``no''; dense and PRA generations begin with ``No'', while truncation begins with
``Yes''. However, PRA routes three tail chunks, none overlapping the annotated evidence, so the
QASPER output cannot be credited to retrieved evidence. Qwen also fails the oracle text-RAG
control on both examples. The proper conclusion is negative: this frozen top-layer,
one-gist-per-chunk configuration has not demonstrated content-causal retrieval.

The selections expose a concrete failure mode. For HotpotQA, evidence lies at token spans
107--137 and 543--588, while routed chunks lie near 1216--1318. For QASPER, evidence lies at
0--317, while selected chunks lie near 4032--4192. Post-RoPE mean-key gists and an unadapted
top-layer query systematically prefer irrelevant late-source chunks in these two probes.

\section{Systems Observations}

The first run reaches 1,345,618,432 allocated GPU bytes, including the complete FP16 model and
temporary activations. Full reference K/V remains on CPU; only selected detail K/V moves to the
GPU. This is evidence that the residency path executes, not a memory-savings comparison: no
matched dense long-context peak is measured here. Likewise, cold source construction, warm
routing, generation, and transfer are reported separately so cache reuse is not hidden.

The current HF path is intentionally eager and unfused. Variable memory rows are padded into a
rectangle before Qwen's native eager function. Flash attention, SDPA, custom kernels, ANN
routing, asynchronous overlap, and serving concurrency remain outside this correctness phase.

\section{Limitations and Next Experiments}

The study has one family, one checkpoint size, one PRA layer placement, one seed for systems
smoke, and two QA examples. It does not compare full dense context where the complete 1,318- or
4,207-token source fits Qwen's advertised window, because the first experiment deliberately
holds the deployment-native limit at 256. It has no matched RAG retriever, no statistical QA
sample, no LoRA adaptation, and no throughput or cost study.

The next experiment should address routing before family breadth. Candidate corrections are:
\begin{enumerate}[leftmargin=*]
\item compare pre-RoPE/content-only and post-RoPE routing gists while keeping stored detail K
post-RoPE;
\item pool hidden-state or normalized pre-position keys for routing without changing native
K/V materialization;
\item increase gist count and PRA layer placement only after evidence recall is measurable;
\item train a small router/gist adapter with span-level evidence supervision while freezing the
pretrained language model;
\item run enough HotpotQA and QASPER examples to estimate retrieval recall, EM, and F1 with
confidence intervals, and include full dense and a real text retriever.
\end{enumerate}

Only after a Qwen configuration retrieves evidence above simple controls should the same adapter
contract be tested on Llama. Llama should be next in the family order; Gemma follows only after
Qwen and Llama no longer require shared-core changes.

\section{Related Work}

PRA differs from text retrieval-augmented generation, which retrieves text and lets the model
encode it in the current prompt~\cite{lewis2020rag}. It retrieves layer-native K/V after a
separate bounded encoding path. This makes PRA complementary to RAG: a text retriever may choose
documents before PRA routes chunks inside cached model state. Sparse-attention systems such as
Longformer and BigBird constrain token-token connectivity inside a forward
pass~\cite{beltagy2020longformer,zaheer2020bigbird}; PRA instead materializes selected state from
an external addressable cache. KV-cache systems emphasize serving capacity and token retention
policies~\cite{kwon2023vllm,zhang2023h2o,li2024snapkv}. Paper~2 asks whether selection can be
inserted beneath an existing pretrained family without changing its disabled behavior.

\section{Conclusion}

One shared PRA core now serves both the controlled model and a thin Qwen adapter. The frozen
Qwen3-0.6B checkpoint passes bit-exact disabled parity, preserves native RoPE and GQA cache
layout, replays selected native K/V through generation, and executes an implicit long-prompt
head under a hard native-operation bound. These are meaningful integration results. They do
not establish useful retrieval. The first unrestricted-QA probes select no annotated evidence,
and even oracle text-RAG is unreliable at this model scale. The Qwen integration is generic and
stable enough to continue routing work, but not yet strong enough to justify multiplying family
adapters. The next milestone is content-causal evidence retrieval on Qwen; Llama follows that
gate.

\bibliographystyle{plain}
\bibliography{../shared/references}
\end{document}
